\chapter{Anophthalmia and Microphthalmia}
\index{Anophthalmia and Microphthalmia}
\label{sec:Anophthalmia and Microphthalmia}

\section{Overview}
\begin{itemize}[noitemsep]
  \item Anophthalmia is the complete absence of an eye, and microphthalmia is an abnormally small eye; both are congenital (present at birth)
  \item They are rare, affecting roughly 1 in 10,000 births, and may affect one or both eyes
  \item They are among the most severe congenital eye defects and require multidisciplinary care
\end{itemize}
\section{Causes}
These conditions happen because of disrupted development of the eye in early pregnancy, before the 8th week of gestation. Causes include genetic mutations, chromosomal abnormalities, infections during pregnancy (toxoplasmosis, rubella, cytomegalovirus), and exposure to toxins or certain medications. Many cases have no identifiable cause.
\section{Risk Factors}
\begin{itemize}[noitemsep]
  \item Family history of eye malformations
  \item Genetic syndromes (e.g., SOX2, OTX2 mutations; chromosomal anomalies)
  \item Maternal infections during early pregnancy (TORCH infections)
  \item Exposure to alcohol, retinoic acid, or certain medications in early pregnancy
  \item Advanced parental age (in some studies)
\end{itemize}
\section{Signs and Symptoms}
\subsection{Common symptoms}
\begin{itemize}[noitemsep]
  \item Small or absent eyeball, visible at birth
  \item Small, underdeveloped eyelids and eye socket (orbit)
  \item Absent or poorly functioning vision in the affected eye
  \item Shallow socket and small palpebral opening
  \item Associated facial, brain, or other systemic anomalies in syndromic cases
\end{itemize}
\subsection{Emergency warning signs}
\begin{itemize}[noitemsep]
  \item Any newborn with suspected eye malformation should be evaluated urgently to confirm the diagnosis and screen for associated anomalies
\end{itemize}
\section{Progression and Stages}
The condition is present at birth and is static in terms of the eyeball itself, but the orbit and facial bones fail to grow normally because the eye's growth stimulus is absent. Without treatment, the face becomes progressively asymmetric. Visual prognosis depends on whether the condition is bilateral and on associated brain or optic nerve abnormalities.
\section{Complications}
\begin{itemize}[noitemsep]
  \item Bilateral severe visual impairment or blindness
  \item Facial asymmetry and underdeveloped eye socket
  \item Prosis (drooping eyelid) and socket contraction
  \item Psychological and social impact on the child and family
  \item Associated systemic malformations (brain, heart, face) in syndromic forms
\end{itemize}
\section{Diagnosis}
\begin{itemize}[noitemsep]
  \item Clinical examination at birth
  \item Ultrasound, CT, or MRI of the eye and brain to assess the eyeball, optic nerve, and brain
  \item Genetic testing to identify an underlying cause
  \item Hearing assessment and evaluation for associated systemic anomalies
  \item Multidisciplinary assessment by an ophthalmologist, geneticist, and pediatrician
\end{itemize}
\section{Prevention}
\begin{itemize}[noitemsep]
  \item Prenatal ultrasound may detect severe cases before birth
  \item Genetic counseling for affected families
  \item Prevention of maternal infections (rubella vaccination before pregnancy)
  \item Avoidance of alcohol and harmful medications during pregnancy
  \item Maternal folic acid supplementation
\end{itemize}
\section{Treatment Options}
\begin{itemize}[noitemsep]
  \item Socket expansion with progressively larger conformers from infancy to promote orbital and facial growth
  \item Fitted prosthetic eye (ocular prosthesis) for cosmetic and developmental reasons
  \item Early referral for visual rehabilitation in bilateral cases
  \item Surgical socket reconstruction and eyelid procedures as the child grows
  \item Coordinated care for any associated systemic conditions
  \item Genetic counseling and support for the family
\end{itemize}
\section{Living with It}
\begin{itemize}[noitemsep]
  \item Begin socket expansion and prosthesis fitting early for best facial outcomes
  \item Work with a multidisciplinary team and early-intervention services
  \item Access visual rehabilitation and mobility training
  \item Seek peer support and counseling for the family
  \item Plan regular follow-up through childhood and adolescence
\end{itemize}
\section{Outlook}
The visual outcome depends on whether the condition is bilateral and on associated brain or optic nerve involvement. With early socket expansion and prosthetics, excellent cosmetic and facial-growth outcomes can be achieved. Children with bilateral involvement benefit greatly from early visual rehabilitation, and many lead full, independent lives.
